% ECONOMICS_PROBLEM: {"id": "E32-331", "title": "Sectoral shock propagation", "jel_code": "E32", "source_id": "OP-0232"}
% !TeX program = xelatex
\documentclass[11pt,a4paper]{article}
\usepackage[margin=23mm]{geometry}
\usepackage{fontspec}
\setmainfont{Latin Modern Roman}
\usepackage{amsmath,amssymb,mathtools}
\usepackage{xcolor,enumitem,booktabs,longtable,array,xurl,hyperref}
\definecolor{muted}{HTML}{53636C}
\hypersetup{colorlinks=true,urlcolor=blue,linkcolor=blue}
\setlength{\parindent}{0pt}
\setlength{\parskip}{5pt}
\newcommand{\R}{\mathbb R}
\newcommand{\E}{\mathbb E}
\newcommand{\N}{\mathbb N}
\newcommand{\ind}{\mathbf 1}
\newcommand{\Var}{\operatorname{Var}}
\newcommand{\Cov}{\operatorname{Cov}}
\newcommand{\supp}{\operatorname{supp}}
\DeclareMathOperator*{\argmin}{arg\,min}
\DeclareMathOperator*{\argmax}{arg\,max}
\newcommand{\entrymeta}[3]{{\sffamily\small\color{muted}\textbf{\texttt{#1}}\quad|\quad #2\par #3\par}}
\newcommand{\Problemref}[1]{\texttt{#1}}
\newcommand{\JELref}[2]{\texttt{#2} (JEL \texttt{#1})}
\begin{document}
\section*{Sectoral shock propagation}
\textbf{Problem ID:} \texttt{E32-331}\quad \textbf{JEL:} \texttt{E32}\par
% BEGIN ECONOMICS STATEMENT
\label{problem:OP-0232}
{\sffamily\small\color{muted}Original subject group: Business cycles and macro dynamics\par}
\label{definitions:problem:30007513}
\textbf{Audit classification:} Background; proposed extension. The cited model already supplies conditional network-amplification results; this broader threshold specification is editorial. Verification concerns the cited version, not first priority or proof that this exact editorial specification remains unsolved on October 8, 2026.
\subsection*{Definitions and precise research target}
\paragraph{Editorial mathematical specification.} Consider a finite economy with sectors \(i=1,\ldots,n\). Sector \(i\) combines labor and intermediate inputs using a constant-returns production function, has productivity \(A_i\), and purchases an expenditure share \(\omega_{ij}\) from sector \(j\). Households have specified preferences and hold debt subject to a collateral constraint whose collateral is a declared function of sectoral profits. Fix the input-share matrix \(\Omega=(\omega_{ij})\), initial balance sheets, substitution elasticities, and a rule for monetary and fiscal policy. Assume an equilibrium exists and specify an equilibrium selection if constraints generate multiplicity.
For a productivity loss \(s>0\) in sector \(j\), let \(Y_j(s)\) be equilibrium real aggregate output and define
\[
R_j(s)=1-\frac{Y_j(s)}{Y_j(0)},\qquad
s_j(q)=\inf\{s>0:R_j(s)\ge q\},
\]
where \(Y_j(0)>0\), and \(q\in(0,1)\) is a predetermined recession threshold and the infimum is \(+\infty\) if the set is empty. Characterize and quantitatively identify how network structure changes \(s_j(q)\) when financial conditions and shock sizes are held fixed. Separate first-order production effects from nonlinear collateral amplification; establish which conclusions survive plausible substitutions and alternative collateral definitions. Validate amplification predictions against independently identified sectoral shocks. Existing sufficient-statistic results for particular models are admissible benchmarks, not evidence that this entire class is unsolved.

Specify the intervention as \(A_j(s)=A_j(0)e^{-s}\), \(s\ge0\), with all other productivities fixed, and evaluate the same selected equilibrium under each intervention. A network comparison is a pair \((\Omega,\Omega')\) with common primitives and a specified rule replacing missing inputs, not a derivative holding endogenous CES expenditure shares artificially fixed. In a Cobb--Douglas benchmark \(y_i=A_i\ell_i^{\alpha_i}\prod_k m_{ik}^{\omega_{ik}}\), impose \(\alpha_i>0\), \(\omega_{ik}\ge0\), and \(\alpha_i+\sum_k\omega_{ik}=1\); specify labor endowment, final demand, and the collateral inequality in real units. The threshold infimum need not be attained or monotone in the network. For multiple equilibria report the thresholds for every admissible selection or a worst-case threshold.
\subsection*{Variants and assumption changes}
The following are \emph{editorial variants}, not additional published open conjectures.
\begin{enumerate}
\item \emph{Unconstrained benchmark.} Remove borrowing constraints but retain the same technologies and final demand; compare \(s_j(q)\) with its constrained counterpart and with the local elasticity \(-\partial\log Y_j/\partial s|_{0}\). Local elasticities cannot determine a finite recession threshold without curvature restrictions.
\item \emph{Endogenous substitution.} Replace fixed Cobb--Douglas coefficients by CES nests with elasticities in a declared compact set. Determine bounds on the recession threshold when relative prices change shares and collateral values simultaneously.
\end{enumerate}
\subsection*{References and source audit}
\begin{itemize}
\item Jorge Miranda-Pinto; Eugenio Rojas; Felipe Saffie; Alvaro Silva. \emph{Connected for Better or Worse? The Role of Production Networks in Financial Crises}. 2026. Locator: Cover and Abstract; Section1, printed pp.1--3; December 2025 version of WP26-1. Primary text checked. \url{https://www.bostonfed.org/-/media/Documents/Workingpapers/PDF/2026/WP2601.pdf}.
\end{itemize}
The cited model already supplies conditional network-amplification results; this broader threshold specification is editorial.
\begin{enumerate}\raggedright
\item\hypertarget{definitions:source:30007513:1}{} Jorge Miranda-Pinto; Eugenio Rojas; Felipe Saffie; Alvaro Silva. 2026. \emph{Connected for Better or Worse? The Role of Production Networks in Financial Crises}. Cover and Abstract; Section1, printed pp.1--3; December2025 version of WP26-1.
\url{https://www.bostonfed.org/-/media/Documents/Workingpapers/PDF/2026/WP2601.pdf}.
DOI: \url{https://doi.org/10.29412/res.wp.2026.01}.
\end{enumerate}

\subsection*{Common conventions: Source mathematical conventions}

These conventions apply to each entry unless explicitly replaced.

\textbf{Models and identification.} A primitive model \(m\) specifies all probability laws, feasible choices, information, equilibrium and measurement rules used by its target. The admissible class \(\mathcal M\) is fixed independently of outcomes. If \(P_O(m)\) is its observable probability law and \(\tau(m)\) the target, its identified set at observable law \(P\) is
\[
 \mathcal I_\tau(P)=\{\tau(m):m\in\mathcal M,\ P_O(m)=P\}.
\]
Point identification means that this set is a singleton. An empty set signals incompatibility of data and assumptions. Coordinate bounds need not describe the joint set. Bounds are sharp only if every retained target is attainable or arbitrarily approachable by compatible models. Finite bounds require explicit restrictions on unobserved quantities; unrestricted confounding, hidden wealth, extrapolation or arbitrary equilibrium selection can yield uninformative sets.

\textbf{Causal interventions.} A potential outcome \(Y_i(p)\) is the outcome under a complete policy regime \(p\), including timing, financing, information and market spillovers. The target population and units are fixed before treatment. With interference, use the complete assignment vector or a justified exposure mapping; individual no-interference notation alone cannot identify an economy-wide intervention. A difference of mean potential outcomes does not require the cross-world joint distribution. Conditioning on post-treatment migration, survival, enrollment or employment changes the estimand and needs separate justification.

\textbf{Statistical statements.} An estimator, confidence set, forecast or test specifies a sampling law, model class and loss/coverage criterion. Uniform \((1-\alpha)\)-coverage means \(\inf_{m\in\mathcal M}\Pr_m\{\tau(m)\in C_n\}\geq1-\alpha\), or an explicitly asymptotic version. The whole parameter space trivially has coverage, so informativeness requires a stated width, risk or power objective. A forecasting improvement is relative to a named benchmark, information set and evaluation population.

\textbf{Equilibrium and optimization.} An equilibrium comprises feasible plans, optimality, beliefs where needed, budget constraints and all market/resource clearing. The primitives or a stated benchmark define each correspondence \(\mathcal E(p;m)\). Nonemptiness is assumed or proved, never inferred just from compact policy sets. A selected-equilibrium target uses a declared selection rule; a robust target quantifies over all permitted equilibria. Use suprema or infima unless attainment is established. An \(\varepsilon\)-optimal policy is feasible and within \(\varepsilon\) of the finite optimum value. Report an empty feasible set separately.

\textbf{Welfare and accounting.} Normative utility, interpersonal normalization, social weights, numeraire, discounting and the use of public receipts are primitives. Behavioral utility need not coincide with the welfare criterion. Household welfare includes ownership income and rents once; adding profits again to compensating variation generally double-counts them. A monetary equivalent is defined by an explicit transfer and price convention, with monotonicity and range conditions. Average effects, mechanism contributions, transfers and welfare losses are different targets. Infinite sums require convergence and dynamic feasible plans require terminal or no-Ponzi conditions.

\textbf{Source versions.} A working-paper date, revision date and journal publication year are distinguished. A DOI for a journal version is not automatically the DOI of an earlier manuscript. Locators refer to the stated version and printed pages where specified. A metadata-only check does not verify a claim about a full-text conclusion. Every entry's source subsection narrows the provenance claim accordingly.
% END ECONOMICS STATEMENT
\end{document}
